%%%%%%%%%%%%%%%%%%%%%%%%%%%%%%%%%%%%%%%
% Cover Letter - Zipline
% Embedded Engineering Intern (Summer 2027)
% Plain XeLaTeX — matches simple resume style (Spectral font)
%%%%%%%%%%%%%%%%%%%%%%%%%%%%%%%%%%%%%%%

\documentclass[11pt]{article}

\usepackage[margin=0.75in, top=0.7in, bottom=0.7in]{geometry}
\usepackage[hidelinks]{hyperref}
\usepackage{parskip}
\usepackage{fontspec}

\setmainfont{Spectral}[
  Path = ../../cv/fonts/,
  UprightFont    = Spectral-Regular.ttf,
  BoldFont       = Spectral-Bold.ttf,
  ItalicFont     = Spectral-Italic.ttf,
  BoldItalicFont = Spectral-BoldItalic.ttf
]

\setlength{\parindent}{0pt}
\setlength{\parskip}{7pt}

% No hyphenation
\hyphenpenalty=10000
\exhyphenpenalty=10000

\pagestyle{empty}
\urlstyle{same}

\begin{document}

% ── HEADER ───────────────────────────────────────────────────────────────────
\begin{center}
  {\LARGE\textbf{Travis Brown}}\\[3pt]
  \href{mailto:browntravis760@gmail.com}{browntravis760@gmail.com} \textbullet\
  518-796-8703 \textbullet\
  \href{https://www.linkedin.com/in/tbrown5932}{linkedin.com/in/tbrown5932}
\end{center}

\vspace{2pt}
\noindent\rule{\linewidth}{0.5pt}
\vspace{10pt}

% ── DATE ─────────────────────────────────────────────────────────────────────
\begin{flushright}
  \today
\end{flushright}

\vspace{4pt}

% ── BODY ─────────────────────────────────────────────────────────────────────
Dear Zipline Recruiting Team,

I am applying for the Embedded Engineering Internship for Summer 2027.
I'm a Computer Science student at RIT in the accelerated combined bachelors and masters
program, with an estimated graduation date of December 2027, and most of what I do outside
of class is embedded firmware. At GE Vernova I built a full I/O subsystem from scratch in
C on an STM32F7 for industrial radios, on the RIT Electric Vehicle Team I write C++
firmware in ThreadX RTOS for a custom motorcycle VCU with a 300\,Hz control loop across
dual CAN buses, and in my free time I am currently building an autonomous fixed-wing flight
controller on an ESP32S3 with FreeRTOS, PID stabilization, and GPS waypoint following,
on a KiCad PCB I designed myself.

The embedded challenges that come with urban autonomous delivery are the kind I find most
interesting: keeping a hovering aircraft precise enough to lower a package to someone's yard,
operating near buildings and people, at a scale that runs continuously worldwide. Those are
the kinds of problems where you have to understand the hardware to understand the firmware.
At GE Vernova, a week before a customer demo, I diagnosed a servo jitter fault by working
through software, RF interference, and PCB layout one by one with engineers from each
discipline before tracing it to a defective servo. That kind of cross-stack debugging is
what I find most engaging.

I've written hard real-time firmware in both ThreadX and FreeRTOS, with a 300\,Hz hard
real-time control loop on the EVT VCU and concurrent sensor-fusion and autonomous-control
tasks on the flight controller. For protocols and interfaces, I have production firmware
experience with CAN/CANopen, SPI, I2C, UART, ADC, PWM, Modbus TCP, and MQTT.

Zipline makes a delivery every 30 seconds. Writing the embedded software running on those
aircraft is exactly the kind of work I want to be doing.

Thank you for your consideration.

\vspace{16pt}
Kind regards,

\vspace{20pt}
Travis Brown

\end{document}
